\begin{abstract}
Monitoring systems store large volumes of floating-point time series, and
most of these data are read only by alerting rules: threshold tests on
window aggregates with a hold duration. A lossy compressor with one error
bound for all samples either wastes space far from the thresholds or
changes alerts near them. We use the alert decisions as the correctness
criterion. For every rule, the reconstructed series must produce the same
sequence of alert states (\textsf{Inactive}, \textsf{Pending},
\textsf{Firing}) as the original, at any evaluation interval and phase.
For PromQL threshold rules over \code{last},
\code{avg/sum/min/max\_over\_time} and \code{rate} we derive a sound error
budget for each sample, computable in linear time. It is the smallest
value, over all windows that contain the sample, of the window's distance
to the threshold divided by a Lipschitz constant. Because the aggregates
are monotone, the budget splits into separate bounds for downward and
upward errors. For Prometheus' extrapolated \code{rate}, an
order-preserving quantizer leaves only the first and last sample of each
window constrained. Our codec combines predictive coding, a range coder and
a \emph{lazy} step choice that keeps the prediction whenever the budget
allows, which suits the asymmetric budget. After encoding, it decodes the
output, re-evaluates every rule and repairs any changed decision, so the
guarantee holds for the floating-point output. Parts of the
floating-point argument are machine-checked with Kani. On 13 real series
from the public Prometheus demo servers (7 and 30 days), the codec keeps
every alert decision (9\,055 firings) at 1.0--1.5 bits per value. This is
4--6.5 times less than a uniform bound tuned on the same data and less
than Compression Safeguards over SZ3 on 19 of 26 rule variants. In an
ablation, the one-sided budget and the lazy step choice together reduce
the size by a factor of 2.4--2.5.
\end{abstract}
